% ch3.tex -- Module 5 ESSENTIALS, Chapter 3: Circular Motion: Angles, Energy, Work and Torque
% Knowledge points only (no questions). Condensed from src/mod5/ch3.tex (same notation,
% same facts). Dot points: M5.IQ2.03 (omega = Delta theta / t), M5.IQ2.04, M5.IQ2.05.
% All diagrams are TikZ copied (some shrunk) from the full chapter: no figure files needed.
% Every number is taken from the full book and checked by
\chapter{Circular Motion: Angles, Energy, Work and Torque}\label{ch3:chap}
\covers{M5.IQ2.04, M5.IQ2.05}

\begin{summarybox}{In this chapter}
\begin{itemize}
  \item Angles in radians, the angular velocity $\omega$ and the speed $v = \omega r$.
  \item Why the centripetal force does no work; how work changes the total energy.
  \item Speeds and forces in a vertical circle.
  \item Torque, $\tau = r_\perp F = rF\sin\theta$, and how torque changes rotation.
\end{itemize}
\end{summarybox}

%% ==================================================================
\section{Angles and angular velocity}\label{ch3:sec:omega}
\covers{M5.IQ2.03}

\begin{bigidea}
Every point on a spinning object turns through the same angle in the same time, so it has
the same angular velocity $\omega$. Points further out move faster.
\end{bigidea}

\begin{keywords}
  \kw{Radian (rad)}{a unit for angles. One full turn is $2\pi$ rad, the same as \ang{360}.}
  \kw{Angular displacement ($\Delta\theta$)}{the angle turned through, measured at the centre.}
  \kw{Angular velocity ($\omega$)}{the angle turned each second. Unit: \si{rad.s^{-1}}.}
\end{keywords}

\begin{sidebyside}[0.33]
  \centering
  \begin{tikzpicture}[scale=0.95]
    \def\R{1.8}
    \fill[hscbrand!6] (0,0) circle (\R);
    \draw[hscmuted!70, line width=0.8pt] (0,0) circle (\R);
    \draw[guide] (0,0) -- (-30:\R);
    \draw[guide] (0,0) -- (40:\R);
    \draw[displacement] (-30:0.85) arc (-30:40:0.85);
    \draw[displacement] (-30:1.7) arc (-30:40:1.7);
    \draw[hscink, line width=0.6pt] (-30:0.35) arc (-30:40:0.35);
    \node[diagramlabel, font=\footnotesize] at (5:0.6) {$\Delta\theta$};
    \node[diagramlabel, text=hscmuted] at (5:1.25) {$s_\mathrm{A}$};
    \node[diagramlabel, text=hscmuted, anchor=west] at (5:1.8) {$s_\mathrm{B}$};
    \coordinate (A) at (40:0.85);
    \coordinate (B) at (40:1.7);
    \draw[velocity] (A) -- ++(130:0.55) node[left, text=hscvel, inner sep=1pt] {$v_\mathrm{A}$};
    \draw[velocity] (B) -- ++(130:1.1) node[left, text=hscvel, inner sep=1pt] {$v_\mathrm{B}$};
    \fill[object] (A) circle (3.5pt);
    \fill[object] (B) circle (3.5pt);
    \node[diagramlabel, anchor=south, inner sep=1pt] at ($(A)+(80:0.2)$) {A};
    \node[diagramlabel, anchor=south, inner sep=1pt] at ($(B)+(80:0.2)$) {B};
    \fill[hscink] (0,0) circle (1.5pt);
  \end{tikzpicture}
  \captionof{figure}{A and B turn through the same angle in the same time.}\label{ch3:fig:sameomega}
\nextside
  \begin{itemize}
    \item One turn $= 2\pi \text{ rad} = \ang{360}$, so $1 \text{ rad} \approx \ang{57.3}$.
    \item To radians: degrees $\times\dfrac{\pi}{180}$, turns $\times\, 2\pi$. \term{rpm} (turns per
          minute) $\times\dfrac{2\pi}{60}$ gives \si{rad.s^{-1}}.
    \item \term{Arc length} (distance along the edge): $s = r\theta$, with $\theta$ in rad (not on
          the sheet).
    \item In \cref{ch3:fig:sameomega}, A and B have the same $\omega$. B is twice as far out, so
          its arc is twice as long and it moves twice as fast.
    \item Every point has the same $\omega$, $T$ and $f$. The \term{tangential speed}
          $v = \omega r$ (speed along the circle) and $a_\mathrm{c} = \omega^2 r$ grow with $r$
          (same $\omega$).
  \end{itemize}
\end{sidebyside}

\begin{formulacard}{Angular velocity}
  \cardformula{\omega = \frac{\Delta\theta}{t}}
  \begin{symbols}
    \sym{$\omega$}{angular velocity: angle turned each second}{\si{rad.s^{-1}}}
    \sym{$\Delta\theta$}{angular displacement: angle turned}{\si{rad}}
    \sym{$t$}{time taken}{\si{s}}
  \end{symbols}
  \whentouse{anything that spins or moves round a circle: wheels, fans.}
  \trap{one turn is $2\pi$ rad, not 1 rad. Change turns or degrees to radians.}
\end{formulacard}

\textbf{From the period, and the speed of a point} \hscbadgenotonsheet
\[ \omega = \frac{2\pi}{T} = 2\pi f \qquad\qquad v = \omega r \qquad\qquad a_\mathrm{c} = \omega^2 r \]
\begin{itemize}
  \item $T$ = period (s), $f$ = frequency (Hz), $r$ = radius (m), not the diameter.
  \item Derive: one turn ($2\pi$ rad) takes one period, so $\omega = \dfrac{2\pi}{T}$. The
        sheet gives $v = \dfrac{2\pi r}{T}$, so ${v = \dfrac{2\pi}{T}\times r = \omega r}$. Put
        $v = \omega r$ into $a_\mathrm{c} = \dfrac{v^2}{r}$: $a_\mathrm{c} = \omega^2 r$.
\end{itemize}

%% ==================================================================
\section{Work and energy in circular motion}\label{ch3:sec:work}
\covers{M5.IQ2.04}

\begin{bigidea}
In uniform circular motion the net force is at \ang{90} to the motion. So it does no work,
and the speed stays constant.
\end{bigidea}

\begin{keywords}
  \kw{Work ($W$)}{energy transferred by a force when the object moves along the direction of
    the force (or of part of it). Unit: joule (J).}
  \kw{Total mechanical energy ($E$)}{kinetic energy plus potential energy: $E = K + U$.}
\end{keywords}

\begin{sidebyside}[0.27]
  \centering
  \begin{tikzpicture}[x=1cm, y=1cm]
    \def\R{1.1}
    \draw[hscmuted!60, line width=0.8pt] (0,0) circle (\R);
    \fill[hscink] (0,0) circle (1.5pt);
    \node[diagramlabel, left, inner sep=3pt, font=\footnotesize] at (0,0) {centre};
    \coordinate (P) at (\R,0);
    \draw[hscink, line width=0.5pt] ($(P)+(-0.22,0)$) -- ($(P)+(-0.22,0.22)$) -- ($(P)+(0,0.22)$);
    \draw[force] (P) -- ++(-0.75,0) node[pos=0.5, below, inner sep=3pt, text=hscforce] {$F_\mathrm{c}$};
    \draw[velocity] (P) -- ++(0,0.95) node[right, inner sep=2pt, text=hscvel] {$v$};
    \draw[displacement] ($(P)+(0.3,0.0)$) -- ++(0,0.45) node[right, inner sep=2pt, text=hscmuted, font=\footnotesize] {$\Delta s$};
    \fill[object] (P) circle (4.5pt);
  \end{tikzpicture}
  \captionof{figure}{The net force is at \ang{90} to the velocity: no work.}\label{ch3:fig:work}
\nextside
  \begin{itemize}
    \item $W = F_\parallel s = Fs\cos\theta$ \hscbadgeonsheet, with $\theta$ the angle between
          the force and the displacement $s$.
    \item The \term{centripetal force} (the net force towards the centre) is at \ang{90} to the
          velocity at every instant (\cref{ch3:fig:work}). So for each tiny step,
          $W = F_\mathrm{c}\,\Delta s\cos\ang{90} = 0$.
    \item \term{Work--energy theorem}: the work done by the net force equals the change in
          kinetic energy, $W_\mathrm{net} = \Delta K$ (not on the sheet). So
          $\Delta K = 0$: $K = \tfrac{1}{2}mv^2$ and the speed stay constant. Only the
          \emph{direction} of the velocity changes.
  \end{itemize}
\end{sidebyside}

\begin{itemize}
  \item If it speeds up, a \emph{tangential} part of the net force (along the path) does work.
  \item \textbf{Horizontal circle or circular orbit:} the height (or orbit radius) does not
        change, so $U$ is constant. $K$ is constant too, so $E = K + U$ is constant.
  \item \textbf{Ferris wheel:} steady speed, so $K$ stays the same, but $U$ grows as a rider
        rises. The seat does positive work on the rider. The energy comes from the motor.
\end{itemize}

\begin{formulacard}[notonsheet]{Work and total energy}
  \cardformula{W_\mathrm{other} = \Delta E = \Delta K + \Delta U}
  \begin{symbols}
    \sym{$W_\mathrm{other}$}{work done by all forces except gravity}{\si{J}}
    \sym{$\Delta E$}{change in total mechanical energy}{\si{J}}
    \sym{$\Delta K$, $\Delta U$}{change in kinetic and in potential energy}{\si{J}}
  \end{symbols}
  \whentouse{a push, a motor or friction does work on the object.}
  \derive{gravity's work is $-\Delta U$, so $W_\mathrm{net} = -\Delta U + W_\mathrm{other} = \Delta K$.
    Add $\Delta U$ to both sides.}
  \trap{do not count gravity twice. Its work is already in $\Delta U$.}
\end{formulacard}

\subsection{Vertical circles}

If only gravity does work (no motor, push or friction), the speed changes round a vertical
circle: \term{non-uniform circular motion} (moving in a circle with a changing speed).

\begin{sidebyside}[0.4]
  \centering
  \begin{tikzpicture}[scale=0.8]
    \def\R{1.7}
    \draw[hscmuted!60, line width=0.8pt] (0,0) circle (\R);
    \draw[hscink!60, line width=0.6pt] (0,0) -- (0,\R);
    \draw[hscink!60, line width=0.6pt] (0,0) -- (\R,0);
    \draw[hscink!60, line width=0.6pt] (0,0) -- (0,-\R);
    \fill[hscink] (0,0) circle (1.5pt);
    \node[diagramlabel, left, inner sep=3pt, font=\footnotesize] at (0,0) {centre};
    % top
    \draw[force] (-0.11,\R) -- (-0.11,1.1) node[left, inner sep=1pt, text=hscforce] {$F_\mathrm{T}$};
    \draw[force] (0.11,\R) -- (0.11,0.85) node[right, inner sep=1pt, text=hscforce] {$mg$};
    \draw[velocity] (0,\R) -- (-0.7,\R) node[left, inner sep=2pt, text=hscvel] {$v$};
    \fill[object] (0,\R) circle (4.5pt);
    \node[diagramlabel, above, inner sep=6pt] at (0,\R) {slowest};
    % side
    \draw[force] (\R,0) -- (0.9,0) node[above, inner sep=2pt, text=hscforce] {$F_\mathrm{T}$};
    \draw[force] (\R,0) -- (\R,-0.85) node[right, inner sep=2pt, text=hscforce] {$mg$};
    \draw[velocity] (\R,0) -- (\R,1.05) node[right, inner sep=2pt, text=hscvel] {$v$};
    \fill[object] (\R,0) circle (4.5pt);
    % bottom
    \draw[force] (0,-\R) -- (0,-0.4) node[right, pos=0.6, inner sep=2pt, text=hscforce] {$F_\mathrm{T}$};
    \draw[force] (0,-\R) -- (0,-2.55) node[right, inner sep=2pt, text=hscforce] {$mg$};
    \draw[velocity] (0,-\R) -- (1.5,-\R) node[right, inner sep=2pt, text=hscvel] {$v$};
    \fill[object] (0,-\R) circle (4.5pt);
    \node[diagramlabel, anchor=north east, inner sep=3pt] at (-0.15,-\R) {fastest};
  \end{tikzpicture}
  \captionof{figure}{A ball on a string in a vertical circle (arrows not to scale).}\label{ch3:fig:vertical}
\nextside
  \begin{itemize}
    \item The tension $F_\mathrm{T}$ points along the string, to the centre, at \ang{90} to the
          motion (\cref{ch3:fig:vertical}). It does no work.
    \item Gravity points down. Going up, part of it slows the ball. Going down, it speeds the
          ball up. So the ball is slowest at the top and fastest at the bottom.
    \item $K + U$ stays the same: the \term{conservation of mechanical energy}.
  \end{itemize}
\end{sidebyside}

\par\noindent\begin{minipage}{\linewidth}
\textbf{Speed: use energy} (no friction):
\begin{center}
  $\displaystyle \tfrac{1}{2}mv_1^2 + mgh_1 = \tfrac{1}{2}mv_2^2 + mgh_2$\\[4pt]
  \hscbadgecombine{$K = \frac{1}{2}mv^2$ and $\Delta U = mg\Delta h$}
\end{center}
Multiply by $\dfrac{2}{m}$: $v_2^2 = v_1^2 + 2g(h_1 - h_2)$. Top to bottom, the drop is the
diameter, $2r$ (not $r$), so $v_\mathrm{bottom}^2 = v_\mathrm{top}^2 + 4gr$.
\end{minipage}
\smallskip

\par\noindent\begin{minipage}{\linewidth}
\textbf{Forces: only at the very top or bottom.} Add the real forces, taking ``towards the
centre'' as positive. Their sum is $F_\mathrm{c} = \dfrac{mv^2}{r}$ (on the sheet), not an extra force.
\begin{center}
  $\displaystyle \text{top: } F_\mathrm{T} + mg = \frac{mv^2}{r} \qquad\qquad
    \text{bottom: } F_\mathrm{T} - mg = \frac{mv^2}{r}$\\[4pt]
  \hscbadgenotonsheet
\end{center}
\end{minipage}
\begin{itemize}
  \item At the top, $mg$ points \emph{to} the centre: $+mg$, not $-mg$.
  \item On a track, the \term{normal force} $F_\mathrm{N}$ (the push of the track or seat) replaces
        $F_\mathrm{T}$. In a \term{dip} (low point of a valley): $F_\mathrm{N} - mg = \dfrac{mv^2}{r}$.
        On a \term{crest} (top of a hill): $mg - F_\mathrm{N} = \dfrac{mv^2}{r}$.
  \item The string goes \emph{slack} (loose) at the top when $F_\mathrm{T} = 0$. Then
        $mg = \dfrac{mv^2}{r}$, so $v = \sqrt{gr}$: the smallest speed at the top. A cart leaves a
        crest if $v > \sqrt{gr}$.
  \item \term{Apparent weight} is the seat's push $F_\mathrm{N}$: above $mg$ in a dip, below $mg$
        on a crest.
\end{itemize}

%% ==================================================================
\section{Torque: what makes things turn}\label{ch3:sec:torque}
\covers{M5.IQ2.05}

\begin{bigidea}
A force can make an object turn about a pivot. Its turning effect is the torque,
$\tau = r_\perp F$. If the torques balance, the rotation does not change.
\end{bigidea}

\begin{keywords}
  \kw{Torque ($\tau$)}{the turning effect of a force. Unit: newton metre (\si{N.m}).}
  \kw{Pivot (axis)}{the point or line that an object turns about, such as a hinge.}
  \kw{Line of action}{the line along which a force acts, drawn longer in both directions.}
  \kw{Perpendicular distance ($r_\perp$)}{the shortest distance, pivot to line of action.}
\end{keywords}

\begin{itemize}
  \item Torque is bigger when the force is bigger, acts further from the pivot (a door handle
        is far from the hinges), and is closer to \ang{90} to the \term{lever} (the bar that turns).
  \item \Cref{ch3:fig:torque} shows two ways to find the torque (same result). $\theta$ is between
        $r$ and $F$. $\sin\ang{90} = 1$ gives the most torque. A force through the pivot gives none.
\end{itemize}

\begin{figure}
  \centering
  \begin{tikzpicture}[x=1cm, y=1cm]
    % ---- (a) perpendicular distance
    \begin{scope}
      \draw[object, rounded corners=1.5pt] (0.25,-0.13) rectangle (3.25,0.13);
      \draw[object] (0:0.34) -- (60:0.34) -- (120:0.34) -- (180:0.34) -- (240:0.34) -- (300:0.34) -- cycle;
      \fill[hscink] (0,0) circle (1.6pt);
      \node[diagramlabel, below, inner sep=5pt] at (0,-0.3) {pivot};
      \coordinate (P) at (3.0,0);
      \coordinate (Q) at (2.25,-1.299);
      \draw[guide] (P) -- (3.9,0);
      \draw[guide] ($(Q)+(240:0.35)$) -- (P);
      \draw[hscbrand, line width=1.1pt] (0,0) -- node[below left, diagramlabel, text=hscbrand, inner sep=1pt] {$r_\perp$} (Q);
      \draw[hscink, line width=0.5pt] ($(Q)+(150:0.2)$) -- ($(Q)+(150:0.2)+(60:0.2)$) -- ($(Q)+(60:0.2)$);
      \draw[hscink, <->, line width=0.5pt] (0,0.42) -- node[above, diagramlabel] {$r$} (3.0,0.42);
      \draw[hscink, line width=0.6pt] ($(P)+(0:0.5)$) arc (0:60:0.5);
      \node[diagramlabel] at ($(P)+(30:0.75)$) {$\theta$};
      \draw[force] (P) -- ++(60:1.3) node[above, inner sep=2pt, text=hscforce] {$F$};
      \fill[hscink] (P) circle (1.4pt);
      \node[diagramlabel] at (1.6,-1.95) {\textbf{(a)} $r_\perp = r\sin\theta$, \ $\tau = r_\perp F$};
    \end{scope}
    % ---- (b) perpendicular force
    \begin{scope}[xshift=6.9cm]
      \draw[object, rounded corners=1.5pt] (0.25,-0.13) rectangle (3.25,0.13);
      \draw[object] (0:0.34) -- (60:0.34) -- (120:0.34) -- (180:0.34) -- (240:0.34) -- (300:0.34) -- cycle;
      \fill[hscink] (0,0) circle (1.6pt);
      \node[diagramlabel, below, inner sep=5pt] at (0,-0.3) {pivot};
      \draw[hscink, <->, line width=0.5pt] (0,0.42) -- node[above, diagramlabel] {$r$} (2.92,0.42);
      \coordinate (P) at (3.0,0);
      \draw[guide] ($(P)+(0.65,0)$) -- ($(P)+(60:1.3)$) -- ($(P)+(0,1.126)$);
      \draw[force, hscforce!40, line width=1pt] (P) -- ++(60:1.3) node[right, inner sep=2pt, text=hscforce!70] {$F$};
      \draw[force] (P) -- ++(0,1.126) node[left, inner sep=5pt, text=hscforce, align=right] {$F\sin\theta$\\(turns)};
      \draw[force] (P) -- ++(0.65,0) node[right, inner sep=3pt, text=hscforce, align=left] {$F\cos\theta$\\(no turning)};
      \draw[hscink, line width=0.6pt] ($(P)+(0:0.35)$) arc (0:60:0.35);
      \node[diagramlabel] at ($(P)+(25:0.58)$) {$\theta$};
      \fill[hscink] (P) circle (1.4pt);
      \node[diagramlabel] at (1.6,-1.95) {\textbf{(b)} $\tau = r \times F\sin\theta$};
    \end{scope}
  \end{tikzpicture}
  \caption{A force $F$ on a spanner: (a) perpendicular distance; (b) perpendicular force.}\label{ch3:fig:torque}
\end{figure}

\begin{formulacard}{Torque}
  \cardformula{\tau = r_\perp F = rF\sin\theta}
  \begin{symbols}
    \sym{$\tau$}{torque: turning effect of the force}{\si{N.m}}
    \sym{$r_\perp$}{perpendicular distance, pivot to line of action}{\si{m}}
    \sym{$F$}{size of the force}{\si{N}}
    \sym{$r$}{distance from the pivot to where the force acts}{\si{m}}
    \sym{$\theta$}{angle between $r$ and $F$ (degree mode)}{\si{\degree}}
  \end{symbols}
  \whentouse{a force turns an object about a pivot: doors, spanners, motors.}
  \trap{measure $r$ from the pivot. If the angle is given from a line at \ang{90} to the lever,
    use $\theta = \ang{90}$ minus that angle. The unit is \si{N.m}, not J.}
\end{formulacard}

\subsection{How torque changes rotation}

A force changes straight-line motion. In the same way, a torque changes turning motion.
\begin{itemize}
  \item Net torque not zero: the object starts to turn, speeds up or slows down. A bigger net
        torque changes the rotation faster.
  \item \term{Balanced} torques (net torque zero): the rotation does not change. A still object
        stays still, and a turning object keeps turning at the same rate. About one pivot:
        $\sum\tau_\mathrm{clockwise} = \sum\tau_\mathrm{anticlockwise}$ \hscbadgenotonsheet
  \item Uses: door handles far from the hinges, long spanners. A bicycle \term{crank} (pedal arm)
        gives the most torque when it is horizontal and the foot pushes straight down.
\end{itemize}

%% ==================================================================
\section{Key experiments}\label{ch3:sec:experiments}
\covers{M5.IQ2.04, M5.IQ2.05}

\begin{experiment}{Investigating torque: $\tau = r_\perp F$}
\aim{To test $\tau = rF\sin\theta$, and to see how an unbalanced torque changes the rotation.}
\method{A metre rule is pivoted at its \term{centre of mass} (where its weight seems to act). A
  spring balance keeps it level against a fixed load, at different $r$ and angles. Then the
  torques are made unbalanced.}
\results{$Fr$ stayed at about \qty{0.59}{N.m}: $F$ against $1/r$ is a straight line through the
  origin. At an angle, $F\sin\theta$ stayed the same. A bigger net torque turned the rule faster.}
\conclusion{Only the part of $F$ at \ang{90} to the rule turns it: $\tau = rF\sin\theta$.}
\significance{It explains door handles, spanners and levers, and motor torque (Module 6).}
\end{experiment}

\begin{experiment}{Work and energy in circular motion}
\aim{To relate the work done on an object moving in a circle to its total energy.}
\method{\textbf{A:} a stopper on a string is whirled in a horizontal circle. A video gives
  $K = \tfrac{1}{2}mv^2$ at 8 points. \textbf{B:} a ball swings down from rest through a height
  $h$. A \term{light gate} (a beam that times how long the ball blocks it) gives its speed.}
\results{\textbf{A:} $K$ stayed at \qty{0.24}{J} all the way round, at the same height.
  \textbf{B:} $K$ at the bottom was 97\,\% of the $U$ lost ($mgh$), at every height.}
\conclusion{\textbf{A:} $E = K + U$ was constant: the tension is at \ang{90} to the motion and
  does no work. \textbf{B:} gravity does the work. The missing 3\,\% is a \term{systematic
  error} (the same in every reading): the light-gate beam is not a thin line, so every speed
  comes out slightly too low. Within this error, the $K$ gained equals the $U$ lost.}
\significance{Energy conservation explains pendulums, rollercoasters and satellites.}
\end{experiment}

%% ==================================================================
\begin{summarybox}{Remember}
\begin{itemize}
  \item One turn is $2\pi$ rad. $\omega = \Delta\theta/t = 2\pi/T$, $v = \omega r$:
        every point has the same $\omega$.
  \item Uniform circular motion: $F_\mathrm{c}$ is at \ang{90} to $v$, so $W = 0$ and $K$ is constant.
  \item Only forces other than gravity change $E$: $W_\mathrm{other} = \Delta E = \Delta K + \Delta U$.
  \item Vertical circle: energy for the speed (top to bottom: drop $2r$). Top:
        $F_\mathrm{T} + mg = mv^2/r$; bottom: $F_\mathrm{T} - mg = mv^2/r$.
  \item $\tau = r_\perp F = rF\sin\theta$, $r$ from the pivot. A net torque changes the rotation.
\end{itemize}
\end{summarybox}
